% backmatter/appendixG.tex —— 附录 G 习题、参考分析与勘误流程
% 纪律：习题必须能用仓内工件作答（spec/ 规范、experiments/ 运行记录、research/cards/ 方法卡）；
% 参考分析不给唯一答案，只给判据与会改变答案的条件。
\chapter{习题、参考分析与勘误流程}
\label{app:exercises}

\begin{failurebox}
一道自进化方向的习题，最常见的失败是“答案是观点”。
例如“这个方法算不算自进化？”如果学生的回答只需要引用某一篇论文的说法，
那这道题考查的是文献检索而不是本书的框架。
本书的习题一律改成可判定的形式：给定一段系统描述，回答它改了 $A_k$ 的哪个分量、
反馈来自验证层级的第几层、用哪个集检验、什么情况下结论会翻。
四问答不上来，说明缺少的不是观点而是信息。
\end{failurebox}

\section{习题的三型}

\begin{itemize}
  \item \textbf{判定型}：给定一段描述或一个真实系统，判断它落在哪一层、属于哪一格、是否构成自进化。
        答错通常是把两层混为一层。
  \item \textbf{计算型}：用给定公式与仓内已有数字算出可复核的量（样本量、区间、回本点表达式）。
        答错通常是重采样单位选错。
  \item \textbf{运行型}：在本书实验台上改一处配置重跑，把结果与原记录对照。
        答错通常是破坏了封存集纪律（例如用封存集挑过候选），
        于是两个运行的结果不再可比。
\end{itemize}

\section{章末题与参考分析的逐题导航}
\label{sec:appG-map}

章末题使用“章号-题号”，G1--G18是另设的补充题集。下表逐题列出相关分析及作答范围；“开放”表示不给唯一标准答案，所列入口与要点帮助检验论证，不能把相关G题当作原题的完整答案。计算行给出给定条件下的结果，仍须说明假设。点击章末题号可回到原题，点击 G 题号可直达相应的展开分析。

\begin{small}
\setlength{\tabcolsep}{3pt}
\begin{longtable}{@{}p{1.1cm}p{4.1cm}p{10.1cm}@{}}
\toprule
题号 & 相关分析入口 & 作答要点与答案范围 \\
\midrule\endfirsthead
\toprule 题号 & 相关分析入口 & 作答要点与答案范围 \\\midrule\endhead
\bottomrule\endfoot

\hyperref[ex:ch01-1]{1-1} & \hyperref[sol:g01]{G1}/\hyperref[sol:g02]{G2}；式\eqref{eq:ak} & 开放判定：把观察到的改变落到具体组件，核对版本和后续加载。只在本次任务里使用的反思可随轨迹归档。 \\
\hyperref[ex:ch01-2]{1-2} & 表\ref{tab:harness}；附录\ref{app:stats} & 开放设计：可描述点估计排序，不能只凭均值区分效应与波动；先指定比较，再补逐题差、重复、选择规则与差异区间。 \\
\hyperref[ex:ch01-3]{1-3} & \hyperref[sol:g04]{G4}；公理1 & 开放设计：评价可变时另保留外部参照，并报告规则变化对分数可比性的影响。 \\
\hyperref[ex:ch01-4]{1-4} & \hyperref[sol:g06]{G6}；第\ref{ch03}章 & 开放设计：固定起点与完整资源，对齐候选选择；把任务内重试同跨任务继承分开。 \\
\hyperref[ex:ch02-1]{2-1} & 表\ref{tab:roles}；式\eqref{eq:delta}；\hyperref[sol:g04]{G4} & 开放判定：列出修改器可编辑的对象。若任务或评分随版本而变，差值仍可计算，但评价系统收益还需共同的外部参照。 \\
\hyperref[ex:ch02-2]{2-2} & 第\ref{sec:ch02-states}节；\hyperref[sol:g05]{G5} & 开放设计：预定任务序列和重复次数，比较逐次复位与共用目录两组。文件检查核对起点，题级通过率差及区间估计效果；即使点差为零，也应报告区间。 \\
\hyperref[ex:ch02-3]{2-3} & 第\ref{sec:ch02-gate}节；式\eqref{eq:gate} & 开放判定：隔离外部写入后，影子版本可并行接收输入，输出不交付；灰度版本会承担部分正式流量。二元门缺少发布前的影子观察阶段。 \\
\hyperref[ex:ch02-4]{2-4} & 表\ref{tab:states}；第\ref{sec:ch02-reset-matrix}节 & 开放设计：按新题清单和选定版本核对起始文件，再查缓存或外部状态。用持久组件的变更与加载记录判断哪里发生了更新。 \\
\hyperref[ex:ch03-1]{3-1} & \hyperref[sol:g06]{G6}/\hyperref[sol:g07]{G7}；表\ref{tab:ch03-verdict} & 开放设计：起始经验、选择机会、完整费用与任务覆盖均可能混淆，逐项设置匹配对照。 \\
\hyperref[ex:ch03-2]{3-2} & 式\eqref{eq:samplesize}、\eqref{eq:ch03-power}；\hyperref[sol:g05]{G5} & 计算：用题设标准差0.29，半宽近似需359题，双侧功效近似需734题，均向上取整。两者都假设独立配对题、稳定方差；后者还用到备择效应和目标功效。 \\
\hyperref[ex:ch03-3]{3-3} & \hyperref[sol:g06]{G6}；第\ref{ch03}章题集划分 & 开放设计：选择过程和独立终评缺一不可；额外候选和封存评价都进入预算。 \\
\hyperref[ex:ch03-4]{3-4} & 第\ref{sec:ch03-control}节；表\ref{tab:ch15-reader} & 开放判定：以同协议的目标主干静态组作对照。格式偏差与答题分可并列观察，两项汇总量无法算出格式偏差造成的失分比例。 \\
\hyperref[ex:ch03-5]{3-5} & 第\ref{ch03}章；\hyperref[sol:g10]{G10} & 开放设计：族内扩题改善条件内估计；新族扩大覆盖。按目标总体决定权重和抽样单位。 \\
\hyperref[ex:ch04-1]{4-1} & 表\ref{tab:ch04-six}；\hyperref[sol:g15]{G15} & 开放设计：列观测来源、可重复条件和不可重放的外部状态，允许多个合理方案。 \\
\hyperref[ex:ch04-2]{4-2} & 第\ref{ch04}章标签与迁移；\hyperref[sol:g07]{G7} & 开放构造：正确但不适用的检索可损害；有噪声的材料也可能保留有用信息，须给具体路径。 \\
\hyperref[ex:ch04-3]{4-3} & \hyperref[sol:g08]{G8}；表\ref{tab:ch04-adj} & 开放设计：臂含义、额外预算、选择规则及反馈来源须先核对，点估计反序不是层级判决。 \\
\hyperref[ex:ch04-4]{4-4} & \hyperref[sol:g05]{G5}/\hyperref[sol:g08]{G8}；附录\ref{app:stats} & 开放分析：扩题、重复与分层针对不同变异；重复与估计方法不增加独立题数。 \\
\hyperref[ex:ch04-5]{4-5} & 第\ref{sec:ch04-framing}节；\hyperref[sol:g15]{G15} & 开放设计：给旧目标之外的可观察结果，如后续重开或最终状态；说明反馈如何进入决策。 \\
\hyperref[ex:ch05-1]{5-1} & 表\ref{tab:ch05-forms}；\hyperref[sol:g18]{G18} & 开放观察：给实际条目、分类依据和可见失败，不以占比推断因果作用。 \\
\hyperref[ex:ch05-2]{5-2} & 式\eqref{eq:ch05-link}；\hyperref[sol:g09]{G9} & 开放反例：共享词义可以掩盖不同适用条件；列出排名及错误采用链。 \\
\hyperref[ex:ch05-3]{5-3} & 第\ref{ch05}章；\hyperref[sol:g16]{G16} & 开放分析：条件文本不等于物理删除、停止读取或撤销已传播内容；可给其他具体反例。 \\
\hyperref[ex:ch05-4]{5-4} & \hyperref[sol:g09]{G9}/\hyperref[sol:g10]{G10}；表\ref{tab:ch15-reader} & 开放分析：内容、来源质量和接触任务范围均可改变；先固定同源材料再解释写法差。 \\
\hyperref[ex:ch06-1]{6-1} & \hyperref[sol:g11]{G11}；表\ref{tab:ch06-who} & 开放判定：分别登记写、留、取、删；新增策略训练、反馈、运行与维护成本。 \\
\hyperref[ex:ch06-2]{6-2} & 式\eqref{eq:ch06-amac}；\hyperref[sol:g11]{G11} & 开放权衡：比较规则可复核性、模型语义判断与费用；不得预定性能必然上升。 \\
\hyperref[ex:ch06-3]{6-3} & 第\ref{ch06}章Mem-$\alpha$；\hyperref[sol:g15]{G15} & 开放设计：评价冲突发现与正确处置，并构造为得分而删除有效经验等投机路径。 \\
\hyperref[ex:ch06-4]{6-4} & 第\ref{ch06}章D-MEM；第\ref{ch13}章 & 开放设计：至少报告触发频率、结果质量与完整资源，并定义请求、尝试、发布的分母。 \\
\hyperref[ex:ch07-1]{7-1} & 第\ref{ch07}章更新流程；\hyperref[sol:g02]{G2} & 开放观察：比较版本差，区分删去、改写与被摘要吸收的旧内容。 \\
\hyperref[ex:ch07-2]{7-2} & 表\ref{tab:ch07-cost}；\hyperref[sol:g14]{G14} & 开放设计：交叉改变重写方式与候选评估调用量，统一任务、停止规则和计时范围。若存在交互，改变一个因素的节省量取决于另一个因素的取值。 \\
\hyperref[ex:ch07-3]{7-3} & 第\ref{ch07}章；\hyperref[sol:g07]{G7} & 开放设计：固定条件下干预条目是否可读，记录实际检索、输入及最终判定。 \\
\hyperref[ex:ch07-4]{7-4} & 表\ref{tab:ch15-reader}；\hyperref[sol:g10]{G10} & 开放设计：冻结同一上下文资产，以新读取者自身静态基线作对应比较。 \\
\hyperref[ex:ch08-1]{8-1} & 第\ref{ch08}章技能元数据；\hyperref[sol:g18]{G18} & 开放观察：先定义误用和暴露次数；库大小与任务难度同时变化时不能直接归因。 \\
\hyperref[ex:ch08-2]{8-2} & 表\ref{tab:ch08-voyager}；\hyperref[sol:g06]{G6} & 已有相同迭代上限下的有无对照；总资源比较还需来源构建、实际尝试及各阶段费用，固定可见反馈与选择机会。 \\
\hyperref[ex:ch08-3]{8-3} & 图\ref{fig:ch08-kstar}；\hyperref[sol:g12]{G12}/\hyperref[sol:g14]{G14} & 按图中假设计算：第4次平价，第5次严格节省。右图第6次加收维护费后技能较贵，第8次再次平价。若从第$m$次加收$M$，自该次起改比较$12+2k+M$与$5k$。 \\
\hyperref[ex:ch08-4]{8-4} & 第\ref{ch08}章；\hyperref[sol:g16]{G16} & 开放设计：预定技能集合、接口前后条件和静默失效标准，保留全部技能分母。 \\
\hyperref[ex:ch09-1]{9-1} & 图\ref{fig:ch09-workflow}；\hyperref[sol:g13]{G13} & 开放计算：直出时计生成求解与评审，转入修正时再加修正调用及等待。误放可能交付错解，误拒可能改坏正确解；净收益还需初答正确率、各分支修正效果和费用。 \\
\hyperref[ex:ch09-2]{9-2} & 图\ref{fig:ch09-workflow}；式\eqref{eq:ch09-latency} & 条件分析：总请求可能增加；时延由调度与队列决定，相关性由生成机制决定，无通用不变量。 \\
\hyperref[ex:ch09-3]{9-3} & 第\ref{ch09}章archive；\hyperref[sol:g18]{G18} & 开放设计：来源版本、评价范围与选择规则是最小起点，另列恢复依赖及跨搜索采用条件。 \\
\hyperref[ex:ch09-4]{9-4} & 第\ref{ch09}章预算；\hyperref[sol:g06]{G6} & 开放设计：固定共同反馈和总资源，预定扁平选择器；预算单位应匹配实际部署约束。 \\
\hyperref[ex:ch10-1]{10-1} & 表\ref{tab:ch10-seal}；\hyperref[sol:g06]{G6} & 开放分析：100是筛选范围内Oracle值；分别固定编辑内容与适配过程才能区分瓶颈。 \\
\hyperref[ex:ch10-2]{10-2} & 表\ref{tab:ch10-seal}、\ref{tab:ch10-cpt} & 单段使用 LoRA，多段使用全参数训练，较大语料档的评价范围也改变。停止规则须预定终评、训练量、真实数据比例和多次查看的处理。 \\
\hyperref[ex:ch10-3]{10-3} & 第\ref{sec:ch10-gate}节；\hyperref[sol:g16]{G16} & 开放判定：记忆条目与适配器分别绑定基座、读取配置及恢复快照；文件恢复不等于效果撤销。 \\
\hyperref[ex:ch10-4]{10-4} & 第\ref{sec:ch10-gradient}节；式\eqref{eq:ch10-normalization} & 开放反例：先说明等价指更新方向还是参数步，再看奖励、步数或筛选条件的改变是否保留原目标。 \\
\hyperref[ex:ch10-5]{10-5} & 第\ref{ch10}章资源；\hyperref[sol:g14]{G14} & 计算：串行且无额外开销时3000--4500秒，即50--75分钟；不能据此推算外部经验产量。 \\
\hyperref[ex:ch11-1]{11-1} & 第\ref{ch11}章章首；表\ref{tab:ch11-filter} & 冻结合成器可能使题目难度与求解器错配，去掉惩罚可能增加重复题。分别跟踪题目难度与改写重复率，保留共同的终评。 \\
\hyperref[ex:ch11-2]{11-2} & 第\ref{sec:ch11-filter}节；\hyperref[sol:g06]{G6} & 开放设计：交叉操纵合成器起点与可更新性，保持求解器、任务范围和资源一致。 \\
\hyperref[ex:ch11-3]{11-3} & 第\ref{sec:ch11-consensus}节 & 计算与设计。二元独立设定下，多数一致率期望为\calcnum{0.623046875}，正确率期望为\calcnum{0.5}。可行性、答案校验与难度分别评价，共同错误需独立依据。 \\
\hyperref[ex:ch11-4]{11-4} & 第\ref{sec:ch11-reward}节；\hyperref[sol:g15]{G15} & 开放定义。预定外部任务目标与可接受差异，再检查行为是否利用评价缺口抬高代理奖励。格式改善本身有价值时，需按该目标判断。 \\
\hyperref[ex:ch11-5]{11-5} & 第\ref{ch11}章；\hyperref[sol:g06]{G6}/\hyperref[sol:g15]{G15} & 开放诊断：选择偏差、训练重复、分布变易都可抬分，分别用封存题、去重和分布记录检查。 \\
\hyperref[ex:ch12-1]{12-1} & 表\ref{tab:ch12-sica}；\hyperref[sol:g04]{G4}/\hyperref[sol:g15]{G15} & 开放审计：代码共处只表明权限需查；实际改评测与反复选候选分别影响测量和独立性。 \\
\hyperref[ex:ch12-2]{12-2} & 第\ref{sec:ch12-archive}节；\hyperref[sol:g13]{G13} & 开放设计：缺编辑能力的高分终端与无效分支都可能入档，按实际选择与功能后代记录区分。 \\
\hyperref[ex:ch12-3]{12-3} & 表\ref{tab:ch12-imp}、\ref{tab:ch12-components} & 开放构造：同时列出起点、选定后代与差值；二者同幅提高时增量不变。修改器贡献还需在匹配起点、预算和组件安排后比较。 \\
\hyperref[ex:ch12-4]{12-4} & 第\ref{sec:ch12-vex}节；\hyperref[sol:g04]{G4} & 开放设计：新旧规则交叉复评或独立外部目标，保留规则版本与桥接费用。 \\
\hyperref[ex:ch12-5]{12-5} & 第\ref{sec:ch12-transfer}节；\hyperref[sol:g17]{G17} & 开放分析：大于零比较方法与自身起点；更可迁移比较方法间相同目标范围下的迁移增量。 \\
\hyperref[ex:ch13-1]{13-1} & 表\ref{tab:ch13-timing}；\hyperref[sol:g14]{G14} & 可把请求数除以任务数、尝试数除以请求数、发布数除以尝试数。在每次请求至多尝试一次、每次尝试至多发布一次时，三项分母均非零，其乘积就是每题发布数。若无请求，仍可报告每题发布数为零，后两项无分母。 \\
\hyperref[ex:ch13-2]{13-2} & \S\ref{sec:ch13-signals}；\hyperref[sol:g11]{G11} & 开放设计：观察信号分布、重复材料与真正新材料的区分，再验证其对目标的预测作用。 \\
\hyperref[ex:ch13-3]{13-3} & 第\ref{sec:ch13-costs}节；\hyperref[sol:g14]{G14} & 开放计算：按更新、验证、存储与检索、推理分别记每轮新增费用。若每轮扫描全部旧条目且各条处理费相同，前$n$轮的扫描总量随$n^2$增长，单轮随$n$增长。 \\
\hyperref[ex:ch13-4]{13-4} & \S\ref{sec:ch13-trace-example}；\hyperref[sol:g15]{G15} & 开放诊断：稳定、阈值过高与信号失灵均相容；补记原始信号、候选和监测运行状态。 \\
\hyperref[ex:ch14-1]{14-1} & \S\ref{sec:ch14-provenance-chain}；\hyperref[sol:g15]{G15} & 开放审计：依次查原标签、察觉事件、写入与读取轨迹；以原执行输出检验不一致。 \\
\hyperref[ex:ch14-2]{14-2} & \S\ref{sec:ch14-rail}；\hyperref[sol:g16]{G16} & 开放分析：哈希检验未改写，不能独自证明事实真实；摘要重算另需来源与转换版本。 \\
\hyperref[ex:ch14-3]{14-3} & 表\ref{tab:ch14-collapse}；\hyperref[sol:g15]{G15} & 开放设计：在固定任务上评价表现，独立审查接纳分布；类别熵可能掩盖类内同质化。 \\
\hyperref[ex:ch14-4]{14-4} & 第\ref{ch14}章框架锁定；\hyperref[sol:g15]{G15} & 开放设计：加入与实际目标相连的独立观察，说明其可见性和修正规则的权限。 \\
\hyperref[ex:ch14-5]{14-5} & 附录\ref{app:stats}；式\ref{eq:ch14-consistency}；\hyperref[sol:g05]{G5} & 开放设计：限定污染、任务和效应方向；精度需半宽，功效还需真值、显著性与错误控制；多次全成功与至少一次成功须分别计算。 \\
\hyperref[ex:ch15-1]{15-1} & 表\ref{tab:ch15-reader}；式\eqref{eq:ch15-reader-interaction} & 计算：两差分别为 $-0.3250$、$+0.1000$，后者是目标读取者内的增量。另在每个读取者内取领域拟合减求解器拟合，再用换读差减同读差，得交互点值 $0.0625$。 \\
\hyperref[ex:ch15-2]{15-2} & 表\ref{tab:ch05-e03a}；\hyperref[sol:g09]{G9} & 审计：拆分A缺少能解释两次判定差异的记录。暂停该静态基线、三种记忆臂相对它的差值与区间，以及依赖它的成本效率解释。三种记忆臂之间仍可直接比较。 \\
\hyperref[ex:ch15-3]{15-3} & 表\ref{tab:ch15-validity}；\hyperref[sol:g10]{G10} & 开放设计：预定共同基线、任务与场景、过程合规检查和成本。样本量须对应主指标及场景内相关性；若主张迁移效果不变，应事先规定等价范围。 \\
\hyperref[ex:ch15-4]{15-4} & 表\ref{tab:appF-cross} & 固定总预算时，少做题会缩小任务覆盖，减少重复会降低逐题估计的精度，缩短输出则会改变允许求解器使用的资源。 \\
\hyperref[ex:ch15-5]{15-5} & 第\ref{ch05}章双字段；\hyperref[sol:g09]{G9} & 开放分析：题级收益点估计反号不证明真实效应反号；双字段的唯一必要性还需替代设计对照。 \\
\hyperref[ex:ch16-1]{16-1} & \hyperref[sol:g05]{G5}；表\ref{tab:ch16-power} & 统计解释：半宽式回答精度问题；功效设计另需零假设、备择效应、显著性水平和目标功效，达到精度题数仍不能修复设计、单位与覆盖。 \\
\hyperref[ex:ch16-2]{16-2} & \hyperref[sol:g17]{G17}；表\ref{tab:ch16-verdicts} & 开放诊断：P1问预算分配，P3问形式与容量交互，P5问时机；分别需要匹配重试、交叉对照和任务流比较。 \\
\hyperref[ex:ch16-3]{16-3} & \hyperref[sol:g09]{G9}/\hyperref[sol:g10]{G10}；第\ref{ch05}章 & 开放设计：形式乘容量，固定来源与检索预算，预定独立族与比较范围。 \\
\hyperref[ex:ch16-4]{16-4} & 第\ref{sec:ch14-claims}节 & $D$ 为旧分减新分。区间含零时方向未定；若范围内某次更新的有效配对比较支持 $D<0$，该例就与普遍退化的主张冲突。 \\
\hyperref[ex:ch16-5]{16-5} & \hyperref[sol:g17]{G17}；表\ref{tab:ch15-reader} & 审计：补齐的是一臂记录；跨日期服务身份、少数读取者与未换harness的范围仍需限定。 \\
\end{longtable}
\end{small}


\section{补充题集G1--G18}

\begin{enumerate}[label=G\arabic*.]

  \item\label{ex:g01} \textbf{（第 1 章，判定型）}
        一个代码助手在每次任务失败后把错误信息追加到一个文本文件，下次任务开始时读最近三条。
        它是否构成自进化？按三个时间尺度分别说明它改了什么。
        \hyperref[sol:g01]{参考分析}，可查：第~\ref{sec:timescale}节区分步、回合和版本，表~\ref{tab:scale}列出各层的观察对象。

  \item\label{ex:g02} \textbf{（第 1 章，判定型）}
        上题的系统若把“最近三条”改成“全部累积”，而改变检索规则本身由另一个模型决定，
        用六问区分修改是怎样产生的、下一任务保留了什么。它会带来哪些新的检验义务？
        \hyperref[sol:g02]{参考分析}，可查：第~\ref{sec:sixq}节的第3问“怎样改”和第5问“保留什么”；再用表~\ref{tab:ch06-who}区分写、留、取、删分别由谁决定。

  \item\label{ex:g03} \textbf{（第 2 章，判定型）}
        把 Updatable 从 $\{P,M\}$ 扩大到 $\{P,M,G\}$，必须新增哪些检验？逐条写出，
        并说明原先对提示和记忆的检验，为什么不足以证明改变工作流也有益。
        \hyperref[sol:g03]{参考分析}，可查：先对照表~\ref{tab:updatable}中 $P$、$M$、$G$ 三行的检验要求，再用表~\ref{tab:roles}沿工作流查出受影响的角色。

  \item\label{ex:g04} \textbf{（第 2 章，判定型）}
        把 $\mathcal V$ 纳入 Updatable 时，为什么单列实验组不够，还必须说明最终报告从哪来？
        \hyperref[sol:g04]{参考分析}，可查：第~\ref{sec:components}节的公理~1，以及第~\ref{ch02}章式~\eqref{eq:delta}对共同评价条件的说明。

  \item\label{ex:g05} \textbf{（第 3 章，计算型）}
        用 E01 的实测配对差标准差 \measured{0.2875} 与区间半宽近似 $n\ge(1.96\,s/h)^{2}$，
        估计半宽目标 $h=\calcnum{0.03}$ 时需要的配对题数，并与本书可用的 \measured{16} 道封存题比较。
        若把半宽目标放宽到 \calcnum{0.10}，需要多少道？这两个选择各自让渡了什么？
        \hyperref[sol:g05]{参考分析}，可查：附录\ref{app:stats}式~\eqref{eq:samplesize}给出区间半宽与题数的关系。完整维护档案在\path{agent/BLOCKERS.md}另存历史判决力缺口。

  \item\label{ex:g06} \textbf{（第 3 章，判定型）}
        一个方法在训练时用过的题上提升明显，在未见题中族混合的检验集上提升较小。
        这两个数的差，有多少种可能的解释？如何设计检验把它们分开？
        \hyperref[sol:g06]{参考分析}，可查：表~\ref{tab:ch03-sets}列出训练、选择验证、回归和封存四类题集各自的用途与纪律。

  \item\label{ex:g07} \textbf{（第 4 章，判定型）}
        E02 实测五个非基线臂的回合通过率点估计全部低于无记忆基线（表~\ref{tab:ch04-e02}）。
        写出三种不互相排除的解释，并说明各自的检验方式。
        \begin{sloppypar}
        \hyperref[sol:g07]{参考分析}，可查：运行记录 \nolinkurl{experiments/runs/E02-20260924T153511Z-421eac/summary.json} 的
        \texttt{gains\_vs\_base} 与 \texttt{sample\_sizing} 两节。
        \end{sloppypar}

  \item\label{ex:g08} \textbf{（第 4 章，计算型）}
        E02 中 \texttt{exec} 减 \texttt{review} 的相邻层级收益差为 \measured{$-0.0125$}，
        这是点估计反序。说明为什么它不足以单独证伪总体层级，并列出 P2b 判决仍缺的设计前提。
        \hyperref[sol:g08]{参考分析}，可查：在表~\ref{tab:ch04-adj}找到相邻差，再看表~\ref{tab:ch16-verdicts}的 P2b 行还要求哪些对照。

  \item\label{ex:g09} \textbf{（第 5 章，判定型）}
        E03拆分A的基线汇总记72/80，保存轨迹记70/80；三个记忆臂的汇总判定分别与各自保存轨迹一致。
        这句话能推出“记忆形式不重要”吗？写出它不能推出的理由与能推出的结论。
        \hyperref[sol:g09]{参考分析}，可查：第~\ref{sec:ch05-e03}节的表~\ref{tab:ch05-e03a}与记录冲突案例说明可以保留哪些直接比较。P3 所需的新对照见表~\ref{tab:ch16-verdicts}。

  \item\label{ex:g10} \textbf{（第 5 章，运行型）}
        为 E03 拆分 A 设计一轮使用另外两个留出族的新运行，先给出离线清单与比较协议。
        要求：记录与旧 \texttt{sortkey} / \texttt{strindex} 划分的差异、完整成本及可比性；付费或执行候选前须具备对应授权与运行隔离。旧轨迹不得覆盖，已完成的 E03c 不在此题重跑范围。
        \begin{sloppypar}
        \hyperref[sol:g10]{参考分析}，可查：\path{experiments/protocols/E03.md}列明原来的两个留出族。按附录\ref{app:repro}的运行标识规则，为新划分另列编号。完整档案中的\path{experiments/lab/run_e03.py}可供核对历史实现。
        \end{sloppypar}

  \item\label{ex:g11} \textbf{（第 6 章，判定型）}
        给一个记忆系统，分别指出准入、组织、触发与操作策略由固定规则、LLM判断还是训练所得策略控制。若将其中一个决策改为学习策略，需要新增哪些对照与资源记录？
        \hyperref[sol:g11]{参考分析}，可查：第~\ref{sec:ch06-object}节区分分数、动作与状态。表~\ref{tab:ch06-who}列出各项决策的归属，第~\ref{sec:sixq}节给出六问的填写示例。

  \item\label{ex:g12} \textbf{（第 8 章，计算型）}
        在固定质量约束和成本口径下，将技能方案开始达到成本持平或严格节省的次数分别定义为阈值；阈值不等于数学上的曲率拐点。
        若在教学设计中只选 \calcnum{1}、\calcnum{3}、\calcnum{5} 三个格点采样，
        为什么不能从有限格点决定整个次数轴上是否存在交点？构造两种可能漏掉交点的曲线，并标明它们是教学设定。
        \hyperref[sol:g12]{参考分析}，可查：第~\ref{sec:ch08-economics}节讨论成本曲线与取样格点，表~\ref{tab:ch16-verdicts}列出 P4 所需的比较。完整维护档案另存原协议 \path{experiments/protocols/E05.md}。

  \item\label{ex:g13} \textbf{（第 9 章，判定型）}
        一个多 Agent 系统修改了角色分工，但每个角色的内部组件都没变。
        按本书定义它是否发生了一次版本变更？登记在哪一个分量里？
        \hyperref[sol:g13]{参考分析}，可查：表~\ref{tab:components}中 $G$ 的定义与第~\ref{sec:ch09-repr}节的工作流表示。

  \item\label{ex:g14} \textbf{（第 13 章，计算型）}
        把回本点写成 $N_{\text{break-even}}=C_{\text{evolve}}/\Delta\text{cost}$。
        说明这个固定成本模型何时成立，零节省或负节省时如何处理；区分连续交点、整数次持平与严格节省，并指出本书真实运行尚缺哪些输入。
        \hyperref[sol:g14]{参考分析}，可查：第~\ref{sec:ch08-economics}节讨论回本条件。核对真实费用时，再看第~\ref{sec:ch13-costs}节的四笔账与附录表~\ref{tab:appF-ledger}的覆盖范围。

  \item\label{ex:g15} \textbf{（第 14 章，判定型）}
        一个系统在运行中自行生成了检验用例，并用自己的验收器评分，成功率持续上升。
        说明还缺少哪些信息才能判断改进是否可信，以及哪些额外证据能够区分自证循环、奖励破解与真实进步。
        \hyperref[sol:g15]{参考分析}，可查：\path{spec/glossary.md}中的“自证循环”和“奖励破解”两条，以及第~\ref{sec:ch14-iso}节的隔离条件。

  \item\label{ex:g16} \textbf{（第 14 章，判定型）}
        回滚之后的系统记为 $A_{k}$。若回滚前曾发生过一次 $L$ 的新增（新增技能进入技能库），
        回滚是否必须把它一起撤掉？写出两种处理各自的检验负担。
        \hyperref[sol:g16]{参考分析}，可查：第~\ref{sec:ch14-rollback}节的分阶段回退，以及表~\ref{tab:ch16-verdicts}中 P6 的退化与恢复问题。

  \item\label{ex:g17} \textbf{（第 16 章，综合型）}
        对照表~\ref{tab:ch16-verdicts}的十个研究问题与表~\ref{tab:appg-historical-claims}的原假设，区分“已有观察 / 下一步对照 / 尚未执行的扩展方向”，
        并对每一条写出：要让它离开当前状态，最少需要补什么。
        \hyperref[sol:g17]{参考分析}，可查：\texttt{spec/propositions.md} 全表。

  \item\label{ex:g18} \textbf{（全书，综合型）}
        任选一篇本书参考文献中的论文，先确认所读版本，用本书的六问把它重建一遍：
        每条结论标明“定位于原文第几节 / 表几”；尚未找到写“未定位”，核对相应范围后仍缺失才写“在所核对范围未报告”，不得补全。
        \hyperref[sol:g18]{参考分析}，可查：公开预览可用\path{research/cards/voyager.md}对照六问各栏，再按所选论文逐项定位。完整档案另附空白模板\path{research/cards/TEMPLATE.md}。

\end{enumerate}

\section{参考分析应该长什么样}

\begin{sloppypar}
一道本书习题的参考分析必须含四部分，缺任一部分都算未完成：
\end{sloppypar}

\begin{enumerate}
  \item \textbf{判据}：回答问题用的是哪一条定义、哪一个公式、哪一次运行的哪个字段。
  \item \textbf{数据与出处}：用到的每个数字标身份（引用 / 计算 / 实测）并能追到文件或运行 ID。
  \item \textbf{会改变答案的条件}：至少写一条“若 X 不成立，则答案变为 Y”。
        没有这一条的参考分析，等于把一次读数写成了定律。
  \item \textbf{未答的部分}：明确列出习题中无法用现有证据回答的子问，并说明缺什么。
\end{enumerate}

\begin{principlebox}
参考分析的目的是让读者能推翻它。
凡是写不出“什么情况下我会错”的答案，都不是分析，是结论。
\end{principlebox}

\section{编号习题的展开分析}
\label{sec:appg-worked}

下面的分析示范如何使用判据、指出缺口和保留会改变结论的条件。
它们不是把开放问题变成固定口令，也不代替读者核对实际系统。

\subsection{G1：读到了旧错误，究竟证明了什么}\label{sol:g01}

题目已经给定两个事实：错误信息在任务后保留，下一任务读取其中一部分。
因此它描述了一个经验驱动的持久状态更新与读取过程，可以定位到 $M$；这比仅把错误打印到日志多了一条生效路径。
但题目没有给出独立检验设计，所以还不能判断更新是否有效，也不能报告任何成功率提升。
回答时把``存在待检验的持久更新''与``更新已被证明有益''分开，才与第1章定义一致。

步层的对象包括模型当下收到的消息与工具反馈，回合记录汇总一次任务的执行过程，版本层对象是下次仍会读取的错误条目及其管理规则。
同一段错误文本可以同时出现在轨迹存档和持久存储中，文件内容相同不意味着角色相同。
要判断它在哪一层发挥作用，应检查谁读取它、读取发生在哪个任务边界之后。

会改变答案的条件至少有三种。
如果所谓``下次''只是同一任务内部的重试，就没有证明跨新任务保留；如果每次启动又覆盖了文件，持久链路没有实际成立；如果读取后的内容没有进入执行器可用状态，只能确认文件被访问，不能确认策略起点改变。
反过来，若新任务确实加载候选并接受冻结验收器检查，就可以研究这种持久更新的效果，哪怕最后测得无益或有害。

最小补充材料不是一句``系统有记忆''，而是一份写入前后清单、一条来源轨迹、一条新任务加载记录和一个冻结旧版本对照。
这些材料能回答对象问题；收益大小还需要更多按协议采集的数据。

\subsection{G2与G3：改变规则，新增了哪一种不确定性}\label{sol:g02}

把最近条目改成全部累积，直接改变了检索或保留策略。
按具体实现，它可能属于持久策略 $P$、记忆管理流程 $G$，或作为 $M$ 的版本化元数据保存；必须先说明组件划分，不能仅凭代码所在目录作答。
如果规则由另一个模型产生，还要记录该模型收到的反馈、候选筛选条件与额外调用成本。
仅有一个模型参与修改，不足以证明修改器本身也在被改进。

新的检验义务包括读取内容是否变长、无关经验是否增加、老任务是否退化，以及规则是否在新任务上保持作用。
如果把全部条目直接拼入上下文，长度预算就成为机制的一部分；如果先检索再裁剪，排序与裁剪规则也需要固定。
不能只检验条目有没有保存，而忽略真正送入求解器的内容。

\leavevmode\phantomsection\label{sol:g03}将 $G$ 纳入可更新集合时，新增的不只是一个可以编辑的文件。
工作流可能改变任务分解、工具调用顺序、角色间的信息传递以及失败后的重试位置。
因此应新增结构合法性检查、端到端执行检查、旧能力回归和总资源核算，同时保留固定旧结构的对照。
原来的提示与记忆结果可以作为背景证据，但不能直接证明新结构的收益。

会改变分析的条件是更新范围。
只修改一个不参与执行的角色名称，与修改决定工具权限的角色配置，影响面不同。
若所谓结构变化仅是日志展示方式改变，便不应把它当成能力机制更新；若实际路由和权限改变，检验必须跟随真正的行为路径扩大。

\subsection{G4：验证器也会变时，如何保留共同标尺}\label{sol:g04}

将验证器更新单列成一组，能够标记它与固定验证器组不同，却不能单独解决结果可比性。
如果每个版本都按自己生成的题目和评分规则报告分数，前后差值混合了系统变化、题目难度变化和评分标准变化。
这个差值可以算出来，但无法直接解释为对同一目标的能力改进。

一种可讨论的设计是，把候选内部可更新的反馈器与最终报告所用的冻结评价分开。
系统可以生成新题帮助自己训练或筛选，但最后选定的版本在没有参与这些更新的共同条件下接受评价。
若研究问题本来就涉及评价目标改变，则应把目标变化本身作为研究对象报告，不能仍沿用固定目标下的收益含义。
本题没有给出足够信息决定哪一种设计更合适，正确答案应保留这个缺口。

\subsection{G5：区间精度计算不能冒充功效保证}\label{sol:g05}

E01 汇总保存的配对差标准差为 \measured{0.28751811537130434}，来源是该运行的样本量计算字段。
把它当作下一阶段规划的方差代入量，使用正态近似的区间半宽公式，可得
\begin{align}
n_{h=0.03}&=\left\lceil(1.96\times0.28751811537130434/0.03)^2\right\rceil
 =\calcnum{353},\\
n_{h=0.10}&=\left\lceil(1.96\times0.28751811537130434/0.10)^2\right\rceil
 =\calcnum{32}.
\end{align}
两者都超过当前 \measured{16} 道封存题；放宽半宽要求降低了规模需求，同时允许更不精确的差异估计。
这里的计算假设、方差估计误差和任务相关性都可能改变所需规模，因此这两个整数不是保证覆盖率的工程承诺。

尤其不能把半宽 $h$ 直接改名为最小可检测效应，随后宣称已经规划了检验功效。
检验功效还涉及备择效应、显著性水平、检验方向、相关结构与多重比较。
本题的有用答案，是复算近似并指出缺少哪些设计量；它不能证明历史 E01 曾按这些整数预注册。
当方差来自已经看过的本轮结果时，应将计算身份写成后验精度规划。

\subsection{G8与G12：不要让一个数字承担两个不同结论}\label{sol:g08}

G8 的相邻点估计反序来自环境结果臂减学习式评审臂。
负差说明本次点估计未按所述方向排序；它不自动证明总体层级相反，更不单凭区间跨零就证伪层级命题。
还需要确认四个实验臂确实代表拟研究的反馈类别、选择了哪些相邻比较，以及怎样控制联合错误。
这些前提尚不完整时，P2b 保持未判定。

\leavevmode\phantomsection\label{sol:g12}G12 可以用两条曲线作答。以下均为同等质量下的教学设定。第 $k$ 次复用后，基准累计成本为 $10k$，技能方案比基准多花 $\Delta(k)$。第一种设定在 $k=1,2,3,4,5$ 时的成本差依次为 $1,0,1,-1,1$，第一次持平在第 2 次，第一次严格节省在第 4 次。第二种设定在前五次的成本差恒为 $1$，从第 6 次起取 $\Delta(k)=6-k$，到第 6 次才持平，第 7 次才严格节省。两种方案的累计成本都逐次增加，在第 1、3、5 次也都比基准多花一个单位，只看这三个格点就分不清这两种情形。要用直线外推实测阈值，还需检验每次复用的成本是否稳定，例如中途是否发生了维修。

此外，成本相等和严格节省不同，调用更少和完成同样工作也不同。
E05 小规模试跑中的动态组正好提供后一种提醒：提前失败减少了调用，却没有建立可复用技能。
所以本题还需要一个完成质量约束，才能将成本比较转成采用策略；缺少质量条件时，只能报告资源账的相交。



\subsection{G6与G9：训练题的优势和形式差异不是同一个效应}\label{sol:g06}

G6至少有三类相容解释：模型或持久物记住了训练题的具体内容；检验题的难度或族比例不同；选择过程偏向了在已见题上表现好的候选。它们都可能造成两组收益不同，仅有两个平均数无法分配各自贡献。可先固定候选，在训练题、同族未见题和新族题上分别保留旧版本对照，使用预定权重计算各组配对差；同时登记哪些题曾用于调参。这能定位差异发生在哪一层，却仍不能让难度完全不同的题组成为纯粹的“是否见过”干预。

若要进一步隔离见题的作用，可在内容和难度可比的任务生成范围内，预先随机分配允许接触与保留条件，并对两种训练过程作共同终评。这里描述的是新设计，没有对应本书已经完成的模型结果。若旧候选已根据检验题结果被选过，则“未见题”这个标签需要重查，不能靠重新命名恢复独立性。

\leavevmode\phantomsection\label{sol:g09}G9先处理记录身份。基线汇总与保存轨迹冲突时，两种算术结果不能择优采用；需要原分析对应的逐次记录、替换理由、评分版本和汇总生成链。现有档案不足以裁定，采用边界统一见第5章\S\ref{sec:ch05-record-conflict}。三个记忆臂的判定与各自保存轨迹一致，可以复算结构化减摘要；它不含基线，不能因另一条证据链冲突而一并撤回。即使有效基线得到恢复，完整方案的差仍混合内容、长度和检索，P3所需的形式与容量交叉比较仍缺。

\subsection{G7：把三个解释分别变成可检查的对象}\label{sol:g07}

第一种解释是经验本身不适用。检查链应从来源任务、保存条目、目标任务的适用条件，直到目标回合实际读取的文本；不能只凭记忆文件存在就判断它干扰了求解。第二种解释是执行条件改变，例如反馈生成或检索占用了原可用于求解的资源。应核对完整调用账与候选选择规则，再决定是否需要等总预算的新比较。第三种解释是有限任务与随机运行导致的估计波动。它要求匹配任务、独立单位和缺失处理的区间或重复设计，而不是把五个负点估计当成五次独立确认。

这些解释并不互斥：不适用的经验可以同时增加费用，少量题也可能放大损害估计。第4章已保存的栈案例可以核对“未写文件却在任务后写出正确经验”的时间关系，不能补回未保存的中间评审全文。答案若用新模型生成一段可能的评审来填空，就把新观察冒充成了历史记录。现有结果不足以在三种解释之间给出确定比例。

\subsection{G10与G11：新划分与学习化都需要明确的干预}\label{sol:g10}

G10的离线交付至少包含新旧划分表、任务内容摘要、来源经验的生成范围、读取规则、共同预算和新运行标识。换掉留出族会改变目标分布；保留随机种子并不能让新旧平均数只相差“泛化能力”。如果候选设计已因旧留出结果被修改，旧结果应保留为探索资料，新划分也要检查是否在其他阶段被接触。运行前先确认资源和隔离条件，运行后按新协议单独报告，不能把两个划分的优点各取一半拼成一个方法。

\leavevmode\phantomsection\label{sol:g11}G11没有指定具体系统，应先确定需要控制的决策对象。固定阈值可以控制准入，模型根据提示生成链接可以参与组织，通过奖励训练的策略可以选择写入或删除等操作；这些对象能够组合，不能排列成必然递进的层级。新增验证应围绕实际机制：训练与终评是否分离、固定策略对照是否使用同样资料、训练和部署费用是否完整、旧内容是否遭到不当修改。若学习只覆盖写入而检索仍固定，结论就限定在写入。调用更大的模型、扩大容量或改换方法名称都不直接构成性能证据。

\subsection{G13与G16：结构改变和恢复对象要落到加载状态}\label{sol:g13}

G13按实际执行路径判断。若角色分工决定谁接收哪项子任务、有哪些权限以及信息如何传递，并被后续任务继承，就改变了$G$；角色内部模型和提示均未变，也不妨碍结构版本发生变化。如果只是把日志里的角色名称改写，执行路由保持原样，则没有因此证明能力机制更新。需要核对的最小材料是前后路由或角色配置、加载记录，以及对应任务轨迹，而不是仅比较组织图的外观。

\leavevmode\phantomsection\label{sol:g16}G16中，若声称精确回到旧$A_k$，新加技能及其索引、引用和缓存必须按旧快照恢复；不能仍加载它却沿用旧版本身份。如果决定保留新技能以避免损失有益功能，得到的是一个混合版本，应重新编号，并针对新旧组件的组合做回归与依赖检查。两条路径都还需检查技能已经造成的外部效果，删掉代码并不会撤销已发送的数据或已完成的操作。题目没有给出技能是否参与故障、是否被其他组件依赖，所以不能预先断言哪一种处置代价更低。

\subsection{G14：连续交点、整数阈值和无法回本}\label{sol:g14}

为避免把全部演化费与对照的零费用混在一起，先记技能方案相对对照的额外固定费为$C$，每次完成同等工作节省的同种资源为$s$。在$C\ge0$、$s>0$且两者保持稳定的模型下，净节省为$Ns-C$，连续交点为$C/s$；非负整数任务数上的“不更贵”阈值是$\lceil C/s\rceil$，严格节省阈值是$\lfloor C/s\rfloor+1$。如果两方案都承担构建费，$C$必须取差额，而不是只抄技能臂的一笔总数。

当$s=0$且$C>0$时，没有有限回本点；$C=s=0$时始终持平，没有严格节省。若$s<0$、$C\ge0$，后续任务不会摊回额外固定费。若额外固定费为负、维护费用跳变或每次质量不同，就超出了上述简式，应回到两条完整成本函数及质量条件重新比较。题目所写的$C_{\text{evolve}}/\Delta\text{cost}$只有在前者确实代表额外固定费、后者为稳定正节省时才具有该解释。

E05 v2小试不足以估出有效技能在稳定复用中的节省，也没有覆盖独立任务族上的质量约束；旧E03换读取者的回合计时和之差同样不能充当该式的单次技能节省。因此可以把条件与代数算清，真实临界次数仍留空。为了理解整数区别，可用本附录后文固定费\calcnum{20}/\calcnum{140}、单次费\calcnum{3}/\calcnum{1}的教学例复算：连续交点与整数持平都是\calcnum{60}次，严格节省从\calcnum{61}次开始；这些是构造数，不是报价或模型测量。

\subsection{G15与G17：有待解释的上升与有待判定的命题}\label{sol:g15}

G15先检查测量对象是否保持一致：任务分布、验收规则、提交选择和资源条件是否改变。再检查评分依据是否可被候选改写，以及外部结果是否支持同方向改善。共同偏差被反复保留是自证循环的候选解释；利用判分或观测漏洞提高分数是奖励破解的候选解释；在独立评价上改善才为真实进步提供相应证据。系统自己生成用例并不单独证明其中任何一种。若新增用例经过独立核查，也可以成为训练或检验材料；关键是用途、独立性和信息路径均有记录。

\leavevmode\phantomsection\label{sol:g17}G17应按研究问题区分观察与设计，原假设和旧状态另见本附录末尾。P8保留为扩展方向，软件检查属于装置证据。P1/P1b首先缺匹配的策略、起始记忆与完整预算比较；P2缺统一执行条件下反馈与更新对象的可识别对照，P2b还缺真实对应四类反馈的操作与联合比较规则。P3缺形式与容量的交叉操纵以及任务族覆盖；P3b需限定读取者、日期版本与迁移范围，已有十臂不解除这些限制。

P4缺有效技能及质量受控的复用成本曲线；P5缺共同任务流与总资源下的触发策略比较；P6须把普遍退化主张和幅度单调性拆开，不能靠一个标准差判定；P7仍需为本书E10明确迁移指标、目标任务族与资源范围。第~\ref{sec:ch12-transferobjects}节所述HyperAgents v2组件交叉对照已在对应数学评分配置中支持元层作用，后续可继续辨认元层内部机制。上述对照分别回答不同问题，不能预设补完某项就会支持原方向。新结果也可能支持相反方向，或仍不足以作结论。研究议程见第16章，历史假设与当时判据按归档阅读。

\subsection{G18：方法卡的空白也是可交付信息}\label{sol:g18}

可用Reflexion作为定位练习：原文第3节将动作、评估和反思分开，Memory段说明反思的保存，Algorithm~1给出trial循环\cite{reflexion-2303}。答卷应把这些定位分别放到六问，而不是把摘要里的效果陈述复制六次。是否跨新任务保存，需要另核对具体实验复位规则；未核对代码时标记“代码未检查”，不能仅凭论文流程图填一个代码提交号。

“未定位”记录当前核验工作到哪一步；“在所核对范围未报告”是范围受限的文献观察；“不存在”则是更强的事实主张。三者不互换。方法卡已有答案也不能代替回看原始来源，尤其是卡片中出现与研究对象无关的数据集、硬件或阶段规模时，应检查是否混入另一篇论文的材料。合格答卷可以保留空白，但不能将空白换成貌似完整的推测。

\section{勘误流程}

\begin{sloppypar}
读者可在\href{https://github.com/Palmer-C/understanding-self-evolving-agents/issues}{公开预览仓库的 Issues} 提交勘误。请列出版次或提交号、章节或表号、原文与核对依据。若涉及复算，附上输入文件名、命令、实际输出与预期差异即可。
维护者在完整档案中按以下五步核对并修订源文件。
\end{sloppypar}

\begin{enumerate}
  \item \textbf{登记}。在 \texttt{agent/BLOCKERS.md} 写一行：日期 / 类型（事实 | 数字 | 口径）/ 位置（文件 + 节号或表号）/ 具体描述。
        不要用“某处措辞”这类无法定位的描述。
  \item \textbf{影响面分析}。列出受影响的章节、命题、表格与已引用该数字的位置。
        数字类勘误必须先查它被哪些二次推导用过（例如样本量反推里用过的标准差）。
  \item \textbf{裁定}。事实类与数字类记录核验者、原始来源及修订依据，不把自动核验称为作者审定；口径类进 \texttt{agent/decision-record.md}，
        写明旧口径、新口径与切换理由。
  \item \textbf{修订与回归}。改源文件后必须重跑 \texttt{book/build.ps1} 与
        \texttt{tools/check\_consistency.py}，两处都要零 error。
  \item \textbf{回写}。同步更新 \texttt{agent/source-map.md}（若涉及来源位置）、
        \texttt{agent/TODO.md} 与本附录受影响习题的参考分析指针。
\end{enumerate}

\begin{pitfallbox}
\textbf{反例：把数字勘误写成“措辞调整”。}
如果把第 5 章某个实测数字改小百分之几而不登记，
那么用它做过分母的所有派生量（配额区间、样本量反推、单位成本）都会静默地跟着变。
所以本书要求：凡涉及 \texttt{\textbackslash measured} / \texttt{\textbackslash cited} / \texttt{\textbackslash calcnum}
的修改，一律按数字类勘误走完整流程，不得合并进“文字润色”的提交。
\end{pitfallbox}

\begin{pitfallbox}
\textbf{反例：把口径变更写进正文却不登记。}
同一套数据换一种重采样单位，区间会变宽或变窄，而正文的表看上去一模一样。
维护者应将这类变更记入完整档案的\path{agent/decision-record.md}，并核对沿用同一口径的其他章节；
否则读者无法知道两次报告为什么不一致。
\end{pitfallbox}

\openquestion{本附录的习题目前只到“可判定”这一层，还没有配套的自动判分。
若要上线自动判分，需要给每道题补一个可执行的检查脚本——
确定性题的检查点可以写成断言；例如，固定持久组件的基线不得发布跨任务更新，但仍可写任务工作文件。
涉及设计取舍的分析题还需评价依据与条件，不能只靠匹配答案字符串判分。
本版保留参考分析，不以自动判分尚未实现推迟已有内容的交付；后续是否建设题库应另行确定范围与验收。}

\section{第 4 章补充练习：从轨迹审计到新设计}

\textbf{栈案例。} 对照第\ref{sec:ch04-stack-case}节，分别列出模型请求、工具实际执行、文件改变与任务后经验。
解释为什么测试可读不等于测试已执行，以及为什么后来写对经验不改变原提交。
\emph{参考分析：}原方法只读取末项；自评臂无写文件动作，验收失败。归档经验属于任务结束后阶段。
原始模型自评全文缺失，不能把经验句当作自评原话。

\textbf{矛盾反馈。} 噪声臂标题翻转而日志保持原样，为什么不足以测量纯标签噪声？
\emph{参考分析：}它同时产生标签错误与线索冲突，模型还可能从日志恢复原判定。
纯标签实验应隔离标签；矛盾线索实验则把两类线索及仲裁方式作为明确操作。
二者均需新协议与新数据，不能用修复后的实现重新解释旧轨迹。

\textbf{迟到结果。} 版本甲的失败在版本乙发布后到达，应记到哪个版本，是否应回退？
\emph{参考分析：}先用运行标识关联甲的任务，保存发生与接收时刻；再检查乙是否仍有同一故障条件。
结果归属与当前处置是两项决定，迟到失败不自动要求回退，重复回调也不增加独立证据数。

\textbf{新的比较。} 为同一批错误提交设计“仅标签”与“标签加诊断”的修复比较。
\emph{参考分析：}固定提交、任务材料、修复者与资源口径，先冻结诊断来源和展示范围，保存原文及展示版本，
用不参与诊断的检查评估修复。它估计局部反馈的作用，不自动证明连续任务上的记忆收益。

\section{第 5 章补充练习：同源不等于同输入}

\textbf{截断。} 同一条栈失败记录的截取版本没有保留出错方法，摘要却保留了修复要点。
设计一种不调用模型的保留策略，并指出它可能损失什么。
\emph{参考分析：}可以优先保留实际改动、相关函数和失败断言，其他材料用来源引用连接。
但定位相关函数本身可能出错，不能把预先知道正确答案的选择器当普通摘要器。

\textbf{检索。} 固定字符二元组检索器，改写经验为什么仍会改变处理条件？
\emph{参考分析：}表示变化改变交集与分母、进而改变所选来源。若研究完整方案，这是机制；
若研究固定内容的格式影响，就需先固定源条目标识。两种问题不要混报。

\textbf{成本。} E03 来源账可以恢复，渲染用量缺失，应当怎样报告？
\emph{参考分析：}区分实测来源用量、按保存条目重建的成功调用数、未知 token 与端到端耗时。
不能把零初始化当零费用，也不能用新运行的价格和用量替代旧运行。

\textbf{撤销。} 一条经验被合并进摘要并缓存到若干执行上下文，仅删原文件是否足够？
\emph{参考分析：}需要检查派生主张、索引与执行上下文版本；保留其他独立证据支持的部分，
对无法确认的自动使用先限制。撤销历史存储与停止当前使用是不同操作。

\section{第 9、13、14 章补充练习：检查器的边界}
\label{sec:appg-checkers}

\textbf{并行下界。} 三个不可中断调用各耗时两个单位，并发上限为二，为什么最长支路与总量约束给出三，实际却至少需要四？
\emph{参考分析：}两份资源不能在第三个时间单位之前完成第三个整块任务；公式是下界，不提供可达调度。增加一个单位的后续聚合后，总时间至少为五。若允许切分或服务条件改变，要重新声明假设。

\textbf{候选与选择。} 独立生成三份候选，每份正确概率为0.6；选择器只能挑选已有答案，在存在正确候选时有0.7概率挑对。应报告什么？
\emph{参考分析。}三次都错的概率为 $0.4^3$，所以至少生成一个正确答案的概率为 \calcnum{0.936}。在这些情形中，选择器有七成机会挑对，交付正确率为 \calcnum{0.6552}。比较两种工作流还要计入选择的费用。若允许选择器改写答案，原来的三个候选即使全错，最终也可能改对；这时还要计算这条路径的概率。

换成两道等权的题。设单次答对的概率分别为 \calcnum{0.3} 和 \calcnum{0.9}，每题独立生成三次。读者能否用两题的平均正确率 \calcnum{0.6} 直接代入刚才的式子？逐题算，覆盖率分别为 \calcnum{0.657} 和 \calcnum{0.999}，平均为 \calcnum{0.828}。先合并再算会得到 \calcnum{0.936}，高估了这两道题的平均覆盖率。

若手里只有已经生成的候选，还要区分真实概率与样本比例。假设一道题留下五份程序，其中两份通过测试。等概率抽取三个不同的候选编号，共有 $\binom{5}{3}=10$ 种取法，只有一种恰好取到三个未通过的程序，覆盖率为 \calcnum{0.9}。若把 $2/5$ 当作已知的单次成功概率再算三次，会得到 \calcnum{0.784}。Chen 等估计 HumanEval 的 $\mathrm{pass}@k$ 时，先对每题的 $n$ 份候选、$c$ 份通过计算 $1-\binom{n-c}{k}/\binom{n}{k}$，再按题平均。这里 $1\le k\le n$，当 $n-c<k$ 时分子为零。当同题候选独立同分布时，题内这一项无偏估计该题生成 $k$ 次至少通过一次的概率；部署时能否挑中通过者，还取决于选择器\cite{chen-humaneval-2107}。

\textbf{新颖性。} E11 v2 的输入新颖性高，是否证明需要写入一条新经验？
\emph{参考分析：}该量只描述任务文本与历史输入的差别，是预定触发信号；它没有判断候选经验是否新颖或正确。第一次输入取一也是控制约定。构建、格式准入与后续有效性分别需要证据，不能由一次触发同时代替。


\textbf{习题\ref{ex:ch13-2}：分数变化与语义区分。}
历史八条记录只有字符相似度，没有逐条语义新信息标签，因此不能由全部放行推断八条都新，也不能由分数集中推断度量毫无区分力。先明确要识别的是文字重复、语义重复还是后续有用性；同一份标注不能自动代表这三个目标。

用一组构造数据练习阈值判断。设六条候选的新颖性分数依次为\calcnum{0.9、0.8、0.7、0.6、0.4、0.2}，相应的“相对于已有记忆含新信息”标签为\calcnum{1、0、1、0、1、0}。这些标签是教学设定，不是对E01文本作出的新判定，也不保证任务收益。若分数达到\calcnum{0.5}就放行，四条被放行，其中两条含新信息，两条不含，另漏掉一条新信息；精确率为\calcnum{2/4}，召回率为\calcnum{2/3}。阈值提高到\calcnum{0.7}时，放行三条，精确率为\calcnum{2/3}，召回率仍为\calcnum{2/3}。分数有方差并不意味着排序与目标标签一致；这两个阈值的比较也没有证明其中一个在新任务上最优。

实际评价需要固定候选与已有记忆版本、明确标注者可见的信息，并记录分歧的处理方法。在开发样本上选择度量和阈值后，再到未参与选择的样本上报告对应标签的精确率、召回率与不确定性。若主张是“值得发布”，还须单独观察发布后的质量与资源；语义新信息本身不能代替这项评价。

\textbf{恢复。} 假设按E12 v2设计进行一次恢复检查：目标文件的字节已经还原，复评得分却低于基线，是否说明文件恢复失败？这是教学情境，不是已取得的E12模型结果。
\emph{参考分析：}文件工件恢复与行为恢复是不同观察。随机响应、服务变化或未纳入快照的状态均可能影响表现。先核对字节，再对齐每道题在更新前、更新后及恢复后的三次判决。即使总分回升，原先失去的题也可能仍然失败，另有题目恰好补足分数（见第~\ref{sec:ch14-recovery-denominator}节）。一次差异还不足以识别原因。

\section{第10—12章补充练习：区分目标、代理和归属}

\subsection{成本交点与可部署性是两个问题}

\textbf{题目。} 第10章教学例中，两种方案的固定费用为\calcnum{20}与\calcnum{140}，每任务费用为\calcnum{3}与\calcnum{1}。先复算交点，再说明为什么到达交点仍不能立即决定部署参数方案。若部署版本在交点之前就被替换，原摊销判断还成立吗？

\emph{参考分析。} 等式$20+3N=140+N$给出$N=60$。这个交点只在所设费用保持稳定、固定费用覆盖范围一致且确有相应复用次数时成立。质量、恢复、延迟和数据适用范围仍需单独满足要求；如果参数方案越过质量允许的下降幅度，费用较低也不符合原决策问题。版本在六十次之前替换，则需要按实际寿命结算，不能用未来不存在的任务摊薄成本。

一个常见错误是把额外回归费放到外部方案，却把参数方案的安全与恢复检查省掉。另一种错误是把外部检索每任务成本视为常数，却没有说明存储规模增长会不会改变它。可接受的答案不必提出唯一费用模型，但必须把会改变交点的假设列出来。此题没有真实价格，答案不应写成任何服务商的美元报价。

\subsection{多数一致为什么不是正确率}

\textbf{题目。} 出题系统对一道题生成十个回答，其中六个相同，另四个也相同但与前六个不同。它得到多数一致率\calcnum{0.6}。构造三个真值情形，说明这一观察能对应哪些正确率；再解释为什么仍可把一致率用作某种难度代理。

\emph{参考分析。} 若六个回答对应真值，正确率为\calcnum{0.6}；若另四个对应真值，为\calcnum{0.4}；若两种回答都错，则为零。只知道一致率无法在三者之间选择。这不意味着一致率毫无用途：它可以作为内部不确定性的信号参与课程选择，但需要单独校准其与正确率、可解性和迁移收益的关系。第11章TTCS中的格式有效性检查也不能替代真值判断。

还应考虑答案规范化。两个等价表达被当作不同字符串，会低估一致性；不同解答被错误归为一类，则会高估它。题目中的计数假定类别已经明确，实际实现还须核对分类规则。把这种预处理成本省略，可能使课程代理看起来比实际更容易获得。

\subsection{最好的曲线没有下降，为什么仍有失败候选}

\textbf{题目。} 候选成绩依次为\calcnum{0.40}、\calcnum{0.55}、\calcnum{0.30}、\calcnum{0.52}。画出逐候选曲线与历史最佳曲线，说明每条曲线可以回答什么。若这四个分数的评价题数不同，比较还缺什么？

\emph{参考分析。} 历史最佳为\calcnum{0.40}、\calcnum{0.55}、\calcnum{0.55}、\calcnum{0.55}，其不下降是最大值操作的性质。逐候选曲线保留失败分支，能提示搜索产物的差异；两者都不能单独说明独立任务上的泛化。若评价范围或样本量不同，还需共同复测或清楚标记分阶段估计，不能将不同精度的分数当成完全相同的测量。

进一步问“哪个修改器更好”时，需要比较多个搜索过程，并把选择预算与候选生成预算计入。一个方法偶然产出最高分，另一个方法稳定产出略低但更便宜的版本，取舍取决于预定研究问题。只挑一条最佳曲线不足以回答所有这些问题。

\subsection{整体迁移不自动定位元层贡献}

\textbf{题目。} 将源任务上选出的任务实现与元实现一并迁到新族，结果优于目标初始系统。设计一项额外比较，帮助区分迁入的执行器与修改器分别贡献了什么，并指出可能的兼容性问题。

\emph{参考分析。} 可以讨论初始/迁入任务实现与初始/迁入元实现的交叉组合，并在共同目标预算下比较。但交叉组合必须可执行；若两者接口绑定，直接拼接产生的失败主要说明兼容性，不能自动解释为元层无迁移价值。需保留适配规则、适配成本和失败分母，也不能根据目标终评结果反复修改接口后再称为零样本迁移。

\section{第14—16章补充练习：从证据到结论}

\subsection{九个回合为什么不是九次技能迁移}

\textbf{题目。} 第15章E05 v2案例有九个主回合、七次通过、两个独特任务和一个任务族。分别说明这些数字对应什么分母，并解释为什么动态臂更少的调用不能直接证明动态修订更省钱。

\emph{参考分析。} 主回合包括不同臂、格点与角色；独特任务和任务族的数量描述覆盖，不能由重复回合扩成更多题。动态格点中来源失败，没有准入技能，后续通过也没有检索到该技能；调用较少的一部分来自提前失败。应先核对质量条件、来源状态和实际采用，再谈资源效率。静态技能被注入后通过也不是因果证明，需要无技能或适当替代条件。

如果把所有失败来源删除再计算“已部署技能的平均费用”，可以得到一个条件化量，但必须连同来源成功概率和失败成本报告。否则该量会把生产可用技能的费用排除在外。此题应使用原始轨迹索引，不要求再次请求模型。

\subsection{零故障记录与零故障概率}

\textbf{题目。} 按第16章的独立同分布伯努利教学模型，三十次均未出现指定故障。写出单侧上限的方程，并列出将结果套到持续更新系统前至少需要检查的两项条件。

\emph{参考分析。} 方程为$(1-p)^{30}=0.05$，上限约\calcnum{9.50\%}。需检查观察是否独立、故障定义是否固定、系统与任务分布是否稳定，以及失败是否完整记录。若系统每轮变化，固定$p$可能不是合适模型；若只是同一道题反复检查，也不能据此覆盖所有任务。答案应保留人为模型身份，不把上限说成实际系统故障率。

\subsection{目标变化以后，哪条历史曲线还能继续画}

\textbf{题目。} 工单Agent将目标从关闭比例改为后续问题解决情况。旧日志没有保存后续结果。能否给旧版本补一个新目标分数并与新版本连接成连续曲线？提出一项可行的替代报告。

\emph{参考分析。} 缺少必要观察时，不能凭旧回复推定后续结果。可以从共同起点重新评价两个版本，也可以从新指标开始建立新序列，并明确前后不可直接比较。若还保存了足够的后续记录，则可以按冻结的新分析规则回算，但应登记这是后续定义的指标。判断依据是可恢复的信息与选择过程，而不是一概禁止或允许事后分析。


\subsection{同一任务的多次尝试，要回答哪种成功问题}

\emph{第\ref{ch14}章习题\ref{ex:ch14-5}参考分析。} 可以把待检验的主张限定为：在指定任务分布、固定执行器和相同完整资源范围下，预先定义的反馈扰动方案，相对同期对照使平均得分下降超过允许幅度$\delta$。$\delta$来自研究目标，不能看完区间后选定。旧E02的两臂各自训练、各自构建经验，因此该方案差不能改称“固定同一来源输出，仅翻转两条标签”的净效应。

取$d_i$为同一任务上对照得分减扰动方案得分。在独立同分布的任务级配对差模型下，规划双侧均值区间半宽$h$时，正态近似为$n\approx(z_{1-\alpha/2}\sigma_d/h)^2$。若要检验前述“下降超过$\delta$”的主张，可设单侧原假设$H_0:\mu_d\le\delta$；预期真实差为$\Delta>\delta$、功效为$1-\beta$时，相应近似是$n\approx[(z_{1-\alpha}+z_{1-\beta})\sigma_d/(\Delta-\delta)]^2$。离阈值的距离$\Delta-\delta$决定所需规模，不能把$h$、$\delta$和$\Delta$互换。先导方差$\sigma_d^2$必须对应同一种处理与比较；任务族相关、多重比较和方差估计误差还会改变规划。增加同一道题的调用次数不等于增加独立任务族。

式~\ref{eq:ch14-consistency}回答的是另一类问题。设同一固定任务运行\calcnum{5}次，其中\calcnum{3}次通过，取$k=\calcnum{2}$；这些整数均为教学设定。从现有记录选取两个不同运行，共有$\binom52=\calcnum{10}$个子集。两个都通过的有$\binom32=\calcnum{3}$个，因而该任务的$\mathrm{pass}^{2}$估计为\calcnum{0.30}；至少一个通过的有$\calcnum{10}-\binom22=\calcnum{9}$个，$\mathrm{pass}@2$估计为\calcnum{0.90}。把样本通过率\calcnum{0.60}直接代入独立试验概率，会得到\calcnum{0.36}和\calcnum{0.84}，并非这里的有限样本组合估计。多任务时先逐任务计算，再按预定规则聚合。若运行间保留学习或环境未复位，便需重新检查固定Agent的独立同分布解释是否成立。

\subsection{目标基线、写法差与读取者交互}

\emph{第\ref{ch15}章习题\ref{ex:ch15-1}参考分析。} 由表~\ref{tab:ch15-reader}，同一个换读取者得分\measured{0.3625}，减源读取者基线\measured{0.6875}得到\calcnum{$-0.3250$}，减目标读取者基线\measured{0.2625}则得到\calcnum{$+0.1000$}。前者描述相对源基线的位置，后者描述相对目标基线的经验增量；换模型前后都应保留本地基线，才能区分这两个比较。

式~\ref{eq:ch15-reader-interaction}比较的是两种写法之差是否随读取者改变。沿用该式的$D,S,C,O$记号，四格通过数依次为$Y_{D,C}$的\measured{34}/\measured{80}、$Y_{S,C}$的\measured{29}/\measured{80}、$Y_{D,O}$和$Y_{S,O}$各\measured{56}/\measured{80}，所以
\[
 I_R=\frac{34-29-56+56}{80}=\calcnum{0.0625}.
\]
同一读取者内若先给两种写法都减去共同静态基线，基线项会在写法相减时抵消。这个代数事实不意味着静态基线失去用途，也不意味着点交互已经确定方向。需要在共同题目上联合保留四格结果，先构造题级交互，再按设计估计不确定性；两条边际区间重叠与否不能替代这一步。现有数据还包含任务内重复、同族题目相关、来源集合不同及跨日期运行，不能把该点值直接解释成纯写法效应。这些计算只使用拆分B，不裁定拆分A的冲突基线。

\subsection{相同任务均分为什么不等于相同场景完成率}

\emph{第\ref{ch15}章习题\ref{ex:ch15-3}参考分析。} 设\calcnum{5}个场景，每个场景\calcnum{2}个任务变体，两个系统都通过其中\calcnum{8}/\calcnum{10}个任务；这些数量是教学设定。两者的TGC都为\calcnum{0.80}。若甲的两次失败分属两个场景，则只有\calcnum{3}/\calcnum{5}个场景全部通过，SGC为\calcnum{0.60}；若乙的两次失败同属一个场景，则\calcnum{4}/\calcnum{5}个场景全部通过，SGC为\calcnum{0.80}。场景内任务的归属会改变场景完成率。这些变体可以有不同要求或初始状态，不是$\tau$-bench中同一固定任务的随机复测；小例子的每场景两题也不是AppWorld原基准的任务数量。

迁移设计应预先固定源、目标的任务与场景清单，来源与评价隔离方式，旧版本及目标静态基线，环境复位条件和完整资源范围，再规定主要指标、聚合及配对单位、失败处理和停止规则。样本量按主要比较所需的精度或功效、先导方差及任务与场景的聚集结构规划，不能把同场景变体或同题重复全部当作独立样本。若“结论不变”指迁移效应保持在预定容忍范围内，应事先给出该范围和相应区间判据；只有对应差异区间完整落入范围，才支持该等价结论。未拒绝零差假设或两条区间重叠，都不等于证明不变。例如会议改期题既要检查指定时间发生了必要变化，也要检查其他会议和参与者等无关对象得到保持；允许随操作改变的更新时间戳则应列入许可变化。若规则还要求事先获得确认，就必须保存确认先于执行的过程证据，正确终态不能单独证明这一条件。选择指标的依据应是要回答的迁移问题，而非哪个数字更高。

\section{历史研究假设与问题编号的对应}
\label{sec:appg-history}

表~\ref{tab:appg-historical-claims}列出原计划的十条方向性假设及对应实验。十条均记为\texttt{open}，表示当时的比较尚未形成支持或反证。原协议和全部数值按原样保留，历次判据见\nolinkurl{spec/propositions.md}。

第16章沿用 P 编号，按研究问题列出已有观察和下一项比较。问题清单见\nolinkurl{spec/research-questions-2026-10-21.md}。

\begin{small}
\begin{longtable}{>{\raggedright\arraybackslash}p{0.7cm}>{\raggedright\arraybackslash}p{7.0cm}>{\raggedright\arraybackslash}p{3.8cm}}
\caption{原研究假设与历史实验索引}\label{tab:appg-historical-claims}\\
\toprule
ID & 原主张（历史计划，非已证结论） & 历史实验索引 \\
\midrule\endfirsthead
\multicolumn{3}{l}{表\thetable（续）}\\
\toprule ID & 原主张（历史计划，非已证结论） & 历史实验索引 \\
\midrule\endhead
\midrule\multicolumn{3}{r}{续下页}\\\endfoot
\bottomrule\endlastfoot
P1 & 等预算下，跨任务持久更新优于任务内重试 & E01；重试与预算 \\
P1b & 跨任务更新优于只读不写的持久基线 & E01；起始经验与只读对照 \\
P2 & 反馈质量的影响不小于更新对象选择的影响 & E01/E02；需统一交叉设计 \\
P2b & 改进幅度沿形式验证、执行反馈、学习式评审、内在信号递减 & E02；反馈操作的对应 \\
P3 & 记忆形式比容量更决定跨任务迁移 & E03；形式与容量 \\
P3b & 拟合到问题域的持久物比拟合到求解器的更跨模型、跨 harness & E03/E03c；拟合目标与读取者 \\
P4 & 技能固化存在临界重复次数，低于它不划算 & E05；技能质量与成本 \\
P5 & 单位资源收益：事件触发优于周期优于每任务，漂移时差距扩大 & E11；更新时机 \\
P6 & 任何持久更新都伴随旧能力退化，且与幅度单调相关 & E12；保持与恢复 \\
P7 & 元层改进难以跨任务族迁移 & E10；元层迁移 \\
\end{longtable}
\end{small}

P8 是十条假设以外的候选，E02b 尚未执行。后续研究须先定设计，再采集新数据。
